\paragraph{Constructing the crossing.}
We selected local data-processing programs from our training corpus whose recorded run succeeded and whose file modifications fell within the associated task's allowed changes. Two programs could form a quartet only if their supporting files agreed wherever their paths overlapped and no supporting file conflicted with either program's source. The construction matched the launch context. Where two single-file programs differed only in the assigned entry-point name, the builder could give them a common entry point after checking that renaming was consistent with the tasks and source. This removes a potential clue: the name in the launcher should not reveal which of the two scopes applies.

For each accepted pair the evaluator required a change made by the first program outside the second task's scope, and a change made by the second outside the first task's scope. It combined the same supporting workspace with each alternative, replayed the programs on that workspace, and kept the observed changes on the evaluation side. The 896 accepted groups used 1,792 program replays. Reusing an observed effect across the two task assignments is appropriate here: the user instruction changes which destinations are authorized, while the bytes of the program and its initial files remain fixed. It would not be appropriate for a program that consumed the varying task as an input.
